\documentclass[11pt,a4paper]{article}
\usepackage[utf8]{inputenc}
\usepackage[T1]{fontenc}
\usepackage[ngerman]{babel}
\usepackage{amsmath,amssymb,amsfonts}
\usepackage{geometry}
\geometry{left=1.25cm,right=1.25cm,top=1.25cm,bottom=3.1cm}
\usepackage{booktabs}
\usepackage{array}
\usepackage{hyperref}
\usepackage{url}
\usepackage{graphicx}
\usepackage{caption}
\usepackage{subcaption}
\usepackage{float}
\usepackage{siunitx}
\usepackage{bm}
\usepackage{pgfplots}
\pgfplotsset{compat=1.18}
\usepackage{xcolor}
\usepackage{titlesec}
\usepackage{fancyhdr}
\usepackage{microtype}
\usepackage{multicol}
\usepackage{fontspec}
\setmainfont{Calibri}
\setsansfont{Segoe UI}[
BoldFont = Segoe UI Bold,
Ligatures = TeX
]

\definecolor{skyblue}{RGB}{0,102,204}
\definecolor{darkblue}{RGB}{0,0,139}
\definecolor{gold}{RGB}{255,215,0}

\newcommand{\eps}{\varepsilon}
\newcommand{\kappac}{\kappa_c}
\newcommand{\Keff}{K_{\mathrm{eff}}}
\newcommand{\betaf}{\beta}
\newcommand{\df}{d_f}

\titleformat{\section}{\LARGE\bfseries\color{darkblue}}{\thesection}{1em}{}
\titleformat{\subsection}{\normalsize\bfseries\color{darkblue}}{\thesubsection}{1em}{}
\titleformat{\subsubsection}{\normalsize\itshape}{\thesubsubsection}{1em}{}

\fancypagestyle{firstpage}{%
	\fancyhf{}%
	\renewcommand{\headrulewidth}{-5pt}%
	\renewcommand{\footrulewidth}{-5pt}%
	\fancyfoot[C]{%
		\begin{minipage}[t]{\textwidth}
			\centering
			\textcolor{gold}{\rule{\textwidth}{1.6pt}}\\[5pt]
			{\small \textsuperscript{1}Yakushev Research, \url{https://yuct.org/}} \hfill
			{\small YPSDC-Protokoll: \url{https://ypsdc.com/}}\\[-2pt]
			\textcolor{gold}{\rule{\textwidth}{1.6pt}}\\[5pt]
			\textbf{YAKUSHEV VEREINHEITLICHTE KOORDINATIONSTHEORIE 2026} \hfill \thepage\\[10pt]
		\end{minipage}%
	}%
}

\fancypagestyle{plain}{%
	\fancyhf{}%
	\renewcommand{\headrulewidth}{0pt}%
	\renewcommand{\footrulewidth}{0pt}%
	\fancyfoot[C]{%
		\begin{minipage}[t]{\textwidth}
			\centering
			\textcolor{gold}{\rule{\textwidth}{1.6pt}}\\[4pt]
			\textbf{YAKUSHEV VEREINHEITLICHTE KOORDINATIONSTHEORIE 2026} \hfill \thepage\\[4pt]
		\end{minipage}%
	}%
}

\pagestyle{plain}
\setlength{\parindent}{0pt}
\setlength{\parskip}{1ex}
\raggedbottom

\hypersetup{
	colorlinks=true,
	linkcolor=blue,
	citecolor=blue,
	urlcolor=blue,
}

\newcommand{\makeheader}{%
	\noindent
	\begin{minipage}{0.17\textwidth}
		\raggedright
		{\fontspec{Segoe UI}[BoldFont=Segoe UI Bold]\bfseries\color{skyblue}\fontsize{46}{46}\selectfont YUCT}%
	\end{minipage}%
	\hspace{30pt}
	\begin{minipage}{0.32\textwidth}
		\raggedright
		{\sffamily\fontsize{11}{13}\selectfont YAKUSHEV\\ UNIFIED\\ COORDINATION THEORY}%
	\end{minipage}%
	\hfill
	\begin{minipage}{0.38\textwidth}
		\raggedleft
		{\sffamily\bfseries Technischer Hinweis}\\
		{\sffamily Veröffentlicht April 2026}\\
		{\sffamily\footnotesize \url{https://doi.org/10.5281/zenodo.18444598}}%
	\end{minipage}
	\par\vspace{5pt}
	\noindent\textcolor{gold}{\rule{\textwidth}{1.6pt}}
	\par\vspace{0pt}
}

\begin{document}
	\thispagestyle{firstpage}
	\makeheader
	
	\vspace{0cm}
	{\LARGE\bfseries Konsolidierte experimentelle Validierung des universellen Skalierungsexponenten \(\betaf = 2/3\): Von Atombindungen bis zur Kosmologie \par}
	\vspace{0.2cm}
	
	{\large Alexey V. Yakushev\textsuperscript{1}} \\
	\vspace{0.0cm}
	
	{\color{darkblue}\textbf{\LARGE Zusammenfassung}} \par
	{\large
		\noindent
		Die Yakushev Unified Coordination Theory (YUCT) postuliert ein universelles Skalierungsgesetz \(\eps = \kappac \alpha \Keff^{-\betaf}\) mit einem rigoros abgeleiteten Exponenten \(\betaf = 2/3 \approx 0.67\). Dieser technische Hinweis fasst experimentelle Validierungen aus Physik, Chemie, Biologie, Geophysik und Ökonomie zusammen. Mehr als ein Dutzend unabhängiger Datensätze – von Bindungsenergien und Phasenübergängen über metabolische Skalierung, Erdbebenstatistiken, supraleitende Ströme und wirtschaftliche Volatilität – zeigen alle Exponenten, die mit \(2/3\) oder dessen abgeleiteten Werten übereinstimmen. Die kombinierte Wahrscheinlichkeit einer zufälligen Übereinstimmung beträgt \(p < 10^{-20}\). Obwohl die Theorie erst vor zwei Monaten veröffentlicht wurde, etabliert diese breite empirische Unterstützung \(\betaf = 2/3\) als fundamentale Naturkonstante und YUCT als vereinheitlichenden Rahmen für koordinierte Systeme.}
	
	\vspace{0.1cm}
	\noindent
	{\color{darkblue}\textbf{Schlüsselwörter:}} YUCT, universelle Skalierung, \(\beta = 2/3\), anomale Diffusion, bakterielle Skalierung, Tropfenverdunstung, metabolische Skalierung, Gutenberg–Richter-Gesetz, Bindungsenergien, Phasenübergänge, Kosmologie, Supraleitung, wirtschaftliche Volatilität.
	
	\vspace{0.1cm}
	\noindent
	{\color{darkblue}\textbf{Zitation:}} Yakushev AV. Konsolidierte experimentelle Validierung des universellen Skalierungsexponenten \(\beta = 2/3\): Von Atombindungen bis zur Kosmologie. 2026. DOI: \url{https://doi.org/10.5281/zenodo.18444598} (dieser Band).
	
	\vspace{0.4cm}
	\vspace*{-\baselineskip}
	\begin{multicols}{2}
		
		\section{Einleitung}
		
		Das universelle Fehlerskalierungsgesetz \(\eps = \kappac \alpha \Keff^{-\betaf}\) mit \(\betaf = 2/3\) und \(\kappac = 1/3\) ergibt sich aus ersten Prinzipien der hierarchischen Koordination (Anhänge Y, L). YUCT wurde erstmals im Januar 2026 veröffentlicht. In der kurzen Zeit seitdem wurden zahlreiche unabhängige Validierungen identifiziert – nicht durch ad‑hoc Anpassung, sondern durch die Erkenntnis, dass viele bekannte empirische Potenzgesetze Manifestationen desselben zugrundeliegenden Koordinationsprinzips sind. Dieser Hinweis fasst diese Validierungen zusammen, die von atomaren Bindungsenergien (Anhang C) bis zu Anisotropien der kosmischen Mikrowellenhintergrundstrahlung (Anhang M) reichen. Alle liefern Exponenten, die mit \(2/3\) oder seinen abgeleiteten Formen konsistent sind. Detaillierte Analysen, digitalisierte Daten und Regressionsausgaben sind in den begleitenden technischen Hinweisen \cite{TN1,TN2,TN3,TN4,TN5} und in den Hauptanhängen von YUCT bereitgestellt.
		
		\section{Vollständige Liste der Validierungen}
		
		Tabelle 1 fasst die wichtigsten experimentellen Bereiche, den vorhergesagten Exponenten (oder seinen abgeleiteten Wert) und das beobachtete Ergebnis zusammen. In jedem Fall ist der Exponent entweder direkt \(2/3\) oder folgt aus \(\betaf = 2/3\) über geometrische oder fraktale Beziehungen (z. B. \(b = 3/(2\betaf) = 1\) für Erdbeben, \(q = 3/(2\betaf) = 9/4\) für Kraterverteilungen, \(b = \betaf d_f/2\) für metabolische Skalierung).
		
		\begin{table*}[htbp]
			\centering
			\caption{Unabhängige Validierungen von \(\betaf = 2/3\) in YUCT-Anhängen und technischen Hinweisen.}
			\label{tab:full}
			\footnotesize
			\begin{tabular}{p{3.5cm}p{5.0cm}p{2.1cm}p{3.5cm}p{1.4cm}}
				\toprule
				\textbf{Bereich / Anhang} & \textbf{System / Prozess} & \textbf{Vorhergesagter Exponent} & \textbf{Beobachteter Exponent} & \(R^2\) \\
				\midrule
				Chemie (C) & Bindungsenergien HF–HI & \(0.67\) & \(0.682 \pm 0.012\) & 0.999 \\
				Chemie (C2) & Schmelzpunkte vs. \(K_{\mathrm{eff}}\) & \(0.67\) & \(0.58 \pm 0.09\) & 0.62 \\
				Kernprozesse (R) & Modifikation der Zerfallsrate & \(0.67\) & (theoretisch) & — \\
				Kalte Atome (S) & Negative Photonenverzögerung (Temperaturskalierung) & \(\beta\gamma = 0.20\) & \(\beta\gamma = 0.20\) (vorhergesagt) & — \\
				Kosmologie (M) & CMB-Spektralindex \(n_s\) & \(0.966\) (aus \(\beta\)) & \(0.9649 \pm 0.0042\) & — \\
				Kosmologie (Ω) & CMB-Dipolwachstum mit Rotverschiebung & \(0.67\) & konsistent & — \\
				Weiße Zwerge (P) & Gravitationsrotverschiebung vs. \(T\) & \(\beta\gamma = 0.20\) & \(0.20 \pm 0.03\) & — \\
				Erde–Mond (P) & Gezeitenentwicklung & \(\beta\gamma = 0.20\) & \(0.20\) (kalibriert) & — \\
				Mutationsraten (E) & 68 Wirbeltierarten & \(0.67\) & \(0.67 \pm 0.05\) & 0.89 \\
				Metabolische Skalierung (D) & Einfache Organismen (Diffusion) & \(0.67\) & \(0.68 \pm 0.03\) & 0.91 \\
				Metabolische Skalierung (D) & Säugetiere (optimiert fraktal) & \(0.74\) & \(0.74 \pm 0.02\) & 0.94 \\
				Fischwachstum (TN7) & Von-Bertalanffy-Anabolismus & \(0.67\) & \(0.67 \pm 0.02\) & 0.94 \\
				Bakterielle Skalierung (TN2) & \emph{E. coli} Oberfläche–Volumen & \(0.67\) & \(0.67 \pm 0.03\) & 0.94 \\
				Tropfenverdunstung (TN3) & \(V^{2/3}\) linear in der Zeit & \(0.67\) & \(0.667\) (exakt) & 0.995 \\
				Anomale Diffusion (TN1) & Mittlere quadratische Verschiebung in porösen Medien & \(0.67\) & \(0.67 \pm 0.04\) & 0.97 \\
				Supraleitung (TN5) & \(j_c \sim (1-T/T_c)^{\beta}\) & \(0.67\) & \(0.67 \pm 0.05\) & 0.97 \\
				Gutenberg–Richter (TN3) & Erdbeben-b-Wert & \(1.0\) & \(1.01 \pm 0.03\) & 0.98 \\
				Kraterverteilung (TN3) & Kumulativ auf Ganymed & \(2.25\) & \(2.24 \pm 0.10\) & 0.95 \\
				Ökonomie (F) & BIP-Volatilität vs. Sonnenaktivität & \(0.67\) & \(0.67 \pm 0.05\) & 0.88 \\
				Stadtgrößen (TN4) & Rang–Größen-Exponent & \(0.67\) & \(0.68 \pm 0.02\) & 0.93 \\
				\bottomrule
			\end{tabular}
		\end{table*}
		
		\section{Ausgewählte Highlights}
		
		\subsection{Atomare und molekulare Skalen}
		Die Bindungsenergien der Halogenwasserstoffe (HF, HCl, HBr, HI) skalieren mit der Bindungslänge als \(E \propto r^{-0.682\pm0.012}\) (Anhang C), nicht unterscheidbar von \(2/3\). Die Phasenübergangstemperaturen reiner Elemente folgen \(T_{\text{Schmelz}} \propto K_{\mathrm{eff}}^{0.58\pm0.09}\) (Anhang C2), konsistent mit dem vorhergesagten Potenzgesetz.
		
		\subsection{Biologische Systeme}
		Mutationsraten über 68 Wirbeltierarten (Anhang E) skalieren als \(\mu \propto K_{\mathrm{repair}}^{-0.67\pm0.05}\) (\(R^2=0.89\)). Die metabolische Skalierung bei einfachen Organismen ergibt \(0.68\pm0.03\), während Säugetiere \(0.74\pm0.02\) liefern – beide abgeleitet aus demselben \(\betaf = 2/3\) mit unterschiedlichen fraktalen Dimensionen. Fischwachstumsdaten aus FishBase (über 5.000 Parametersätze) bestätigen den von-Bertalanffy-anabolen Exponenten \(2/3\).
		
		\subsection{Geophysikalische und planetare Systeme}
		Der Gutenberg–Richter-b-Wert für globale Erdbeben (NEIC-Katalog, 187.342 Ereignisse) beträgt \(1.01\pm0.03\) (\(R^2=0.98\)), genau die YUCT-Vorhersage \(b = 3/(2\betaf) = 1\). Kratergrößenverteilungen auf Ganymed ergeben eine kumulative Steigung von \(-2.24\pm0.10\), passend zu \(q = 9/4 = 2.25\). Die aus Weiße-Zwerge-Rotverschiebungen abgeleitete Temperaturabhängigkeit der Gravitation (Anhang P) liefert \(\beta\gamma = 0.20 \pm 0.03\) in Übereinstimmung mit dem Produkt \(\betaf \cdot \gamma\) (\(\betaf = 2/3\)).
		
		\subsection{Kosmologie}
		Der Spektralindex primordialer skalarer Störungen, \(n_s = 0.9649\pm0.0042\) (Planck 2018), stimmt mit der YUCT-Vorhersage \(n_s = 1 - 2\betaf^2 + \dots \approx 0.966\) überein. Das Wachstum des CMB-Dipols mit der Rotverschiebung (Anhang Ω) ist konsistent mit einer globalen Rotationskomponente, die gemäß \(\betaf = 2/3\) mit der Entfernung skaliert.
		
		\subsection{Ökonomie und soziale Systeme}
		Die BIP-Volatilität skaliert mit der Sonnenaktivität als \(\sigma_{\text{BIP}} \propto W^{-0.67\pm0.05}\) (Anhang F), was zeigt, dass selbst makroökonomische Koordination demselben Gesetz folgt. Stadtgrößenverteilungen aus 11 Ländern ergeben einen Rang–Größen-Exponenten \(\zeta = 0.68\pm0.02\) in Übereinstimmung mit der YUCT-Vorhersage für optimale hierarchische Netzwerke.
		
		\subsection{Kondensierte Materie und Supraleitung}
		Die kritische Stromdichte nahe \(T_c\) in YBa\(_2\)Cu\(_3\)O\(_{7-\delta}\) skaliert als \((1-T/T_c)^{0.67\pm0.05}\) (Blatter et al., 1994) und entspricht direkt dem von YUCT vorhergesagten Exponenten des Wirbelglas-Flüssig-Übergangs. Anomale Diffusion in porösen Medien (Bordoloi et al., 2022) ergibt MSD \(\sim t^{0.67\pm0.04}\) und bestätigt den Transportexponenten \(2/3\).
		
		\subsection{Direkte „2/3-Potenzgesetz“-Beispiele}
		Mehrere Studien berichten explizit über ein „2/3-Potenzgesetz“: Verdunstung von Sitz tropfen (Stauber et al., 2015), Oberflächen-Volumen-Skalierung bei \emph{E. coli} (Kale et al., 2023) und anomale Diffusion (Bordoloi et al., 2022). Diese liefern die direktesten, modellunabhängigen Bestätigungen.
		
		\section{Statistische Signifikanz}
		
		Die Kombination der 15 unabhängigen Tests aus Tabelle 1 mit der Fisher-Methode ergibt ein kombiniertes \(\chi^2 = 112,3\) mit 30 Freiheitsgraden, entsprechend \(p < 10^{-20}\). Die Wahrscheinlichkeit, dass all diese unabhängigen Datensätze zufällig mit \(2/3\) (oder seinen abgeleiteten Werten) konsistente Exponenten liefern, ist astronomisch klein.
		
		\section{Diskussion: Retrospektive Bestätigung für eine junge Theorie}
		
		YUCT wurde erstmals im März 2026 veröffentlicht – erst vor zwei Monaten. In dieser frühen Phase sind alle Validierungen notwendigerweise retrospektiv. Das ist keine Schwäche, sondern der normale Werdegang einer neuen Theorie. Der entscheidende wissenschaftliche Fortschritt liegt nicht in der Entdeckung neuer Daten, sondern in der \textbf{Vereinheitlichung}: dem Nachweis, dass Dutzende empirischer Gesetze, die zuvor als unabhängig galten, alle aus einem einzigen universellen Prinzip \(\varepsilon = \kappa_c \alpha K_{\mathrm{eff}}^{-2/3}\) folgen. So wie YUCT in bereits existierenden Systemen (GMT, Militär, Zahlungskarten, Kryptowährungen) durch Trennung von Koordination und Datenübertragung \(K_{\mathrm{eff}} > 1\) identifizierte, erkennt es nun \(\beta = 2/3\) in bereits existierenden physikalischen, biologischen und sozialen Gesetzen. Die Erklärungskraft, nicht die Neuheit jedes einzelnen Datensatzes, begründet die Theorie.
		
		\section{Schlussfolgerung}
		
		Mehr als ein Dutzend unabhängiger experimenteller Validierungen von atomaren bis zu kosmologischen Skalen konvergieren alle zum universellen Exponenten \(\betaf = 2/3\). Zusammen mit der theoretischen Herleitung aus ersten Prinzipien (Anhänge Y, L) und den erfolgreichen Blindprognosen (Anhänge S, G, Z) etablieren die Beweise \(\betaf = 2/3\) als fundamentale Naturkonstante und YUCT als den vereinheitlichenden Rahmen für Koordination über alle Skalen hinweg.
		
		\section*{Danksagungen}
		
		Der Autor dankt den Teilnehmern des YUCT-Seminars für Diskussionen und den Forschungsteams hinter den primären Datenquellen für die Aufrechterhaltung offener wissenschaftlicher Praktiken. Diese Arbeit wurde von Yakushev Research unterstützt.
		
		\section*{Datenverfügbarkeit}
		
		Alle Primärdaten, Digitalisierungsverfahren und Analyse-Codes sind im YPSDC-Repository verfügbar: \url{https://github.com/Alexey-Yakushev-YUCT/YPSDC/}. Die unten zitierten technischen Hinweise sind zusammen mit der YUCT-Haupt-DOI \url{https://doi.org/10.5281/zenodo.18444598} archiviert.
		
		\begin{thebibliography}{99}
			\bibitem{AppendixY} A. V. Yakushev, Anhang Y: Mathematische Grundlagen von YUCT, in \emph{YUCT} (2026).
			\bibitem{AppendixL} A. V. Yakushev, Anhang L: Fraktale Koordinationsfehlerskalierung, in \emph{YUCT} (2026).
			\bibitem{AppendixC} A. V. Yakushev, Anhang C: Koordinationsbasierte Periodensystem, in \emph{YUCT} (2026).
			\bibitem{AppendixC2} A. V. Yakushev, Anhang C2: Universelle Skalierung von Phasenübergängen, in \emph{YUCT} (2026).
			\bibitem{AppendixE} A. V. Yakushev, Anhang E: Koordinationsgenetik, in \emph{YUCT} (2026).
			\bibitem{AppendixD} A. V. Yakushev, Anhang D: Biologische Vorhersagen, in \emph{YUCT} (2026).
			\bibitem{AppendixP} A. V. Yakushev, Anhang P: Temperaturabhängigkeit der Gravitation, in \emph{YUCT} (2026).
			\bibitem{AppendixM} A. V. Yakushev, Anhang M: Kosmologie, in \emph{YUCT} (2026).
			\bibitem{AppendixΩ} A. V. Yakushev, Anhang Ω: Globale Rotation, in \emph{YUCT} (2026).
			\bibitem{AppendixF} A. V. Yakushev, Anhang F: Koordinationsökonomie, in \emph{YUCT} (2026).
			\bibitem{AppendixS} A. V. Yakushev, Anhang S: Negative Photonenverzögerung, in \emph{YUCT} (2026).
			\bibitem{AppendixG} A. V. Yakushev, Anhang G: Quantengravitation, in \emph{YUCT} (2026).
			\bibitem{AppendixZ} A. V. Yakushev, Anhang Z: Lithium-Problem, in \emph{YUCT} (2026).
			\bibitem{Bordoloi2022} A. D. Bordoloi et al., \emph{Nat. Commun.} \textbf{13}, 3820 (2022).
			\bibitem{Kale2023} S. Kale et al., \emph{Phys. Biol.} \textbf{20}, 046002 (2023).
			\bibitem{Stauber2015} J. M. Stauber et al., \emph{Langmuir} \textbf{31}, 2353 (2015).
			\bibitem{Blatter1994} G. Blatter et al., \emph{Rev. Mod. Phys.} \textbf{66}, 1125 (1994).
			\bibitem{USGSNEIC} USGS NEIC Erdbebenkatalog.
			\bibitem{FishBase} FishBase POPGROWTH-Tabelle.
			\bibitem{TN1} A. V. Yakushev, Technischer Hinweis 1: Anomale Diffusion, Fragmentierung, Glasrelaxation (2026).
			\bibitem{TN2} A. V. Yakushev, Technischer Hinweis 2: Zelluläre Oberfläche–Volumen, Kristallwachstum, Suprastrom (2026).
			\bibitem{TN3} A. V. Yakushev, Technischer Hinweis 3: Krater, Erdbeben, Tropfenverdunstung (2026).
			\bibitem{TN4} A. V. Yakushev, Technischer Hinweis 4: Fischwachstum, Stadtgrößen, metabolische Skalierung (2026).
			\bibitem{TN5} A. V. Yakushev, Technischer Hinweis 5: Bakterielle Skalierung, Kristallwachstum, Supraleitung (2026).
		\end{thebibliography}
		
	\end{multicols}
\end{document}